% Generated draft: requires mathematical and source review.
\documentclass[11pt]{article}
\usepackage[margin=25mm]{geometry}
\usepackage{fontspec}
\setmainfont{Latin Modern Roman}
\usepackage{amsmath,amssymb,mathtools,unicode-math}
\setmathfont{Latin Modern Math}
\usepackage{enumitem,needspace,xurl,hyperref}
\setlength{\parindent}{0pt}
\setlength{\parskip}{5pt}
\setlength{\emergencystretch}{3em}
\newcommand{\N}{\mathbb{N}}
\newcommand{\Z}{\mathbb{Z}}
\newcommand{\Q}{\mathbb{Q}}
\newcommand{\R}{\mathbb{R}}
\newcommand{\C}{\mathbb{C}}
\newcommand{\F}{\mathbb{F}}
\newcommand{\A}{\mathbb{A}}
\newcommand{\PP}{\mathbb{P}}
\newcommand{\E}{\mathbb{E}}
\newcommand{\eps}{\varepsilon}
\newcommand{\dd}{\,\mathrm{d}}
\begin{document}
\section*{UnsolvedMath ID: 801027}
\textbf{Source problem code:} catalogue entry
\subsection{Definitions and mathematical statement}
Let \(G\leq\operatorname{Sym}(X)\) be a permutation group. Its action is called \emph{oligomorphic} if, for every integer \(n\geq 1\), the coordinatewise action
\[
 g\cdot(x_1,\ldots,x_n)=(gx_1,\ldots,gx_n)
\]
has finitely many orbits on \(X^n\). In the usual countable setting, \(X\) is countably infinite and \(G\) carries the topology of pointwise convergence, inherited from \(\operatorname{Sym}(X)\) with \(X\) discrete. The supplied statement does not specify the underlying set, whether this topological setting is intended, or whether \(G\) must be closed in \(\operatorname{Sym}(X)\).

For an abstract group \(G\), let \(\operatorname{Aut}_{\mathrm{abs}}(G)\) be its group of abstract automorphisms and define
\[
\operatorname{Inn}(G)=\{c_g:g\in G\},\qquad c_g(h)=ghg^{-1},\qquad
\operatorname{Out}_{\mathrm{abs}}(G)=\operatorname{Aut}_{\mathrm{abs}}(G)/\operatorname{Inn}(G).
\]
For a topological group \(G\), another possible interpretation is
\[
\operatorname{Out}_{\mathrm{top}}(G)
=\operatorname{Aut}_{\mathrm{top}}(G)/\operatorname{Inn}(G),
\]
where \(\operatorname{Aut}_{\mathrm{top}}(G)\) consists of group automorphisms that are homeomorphisms. These interpretations are distinguished here because the source does not specify which is intended.

A topological group \(H\) is \emph{profinite} if it is isomorphic as a topological group to an inverse limit of finite discrete groups.

\textbf{Question (supplied Problem 39).} Is \(\operatorname{Out}(G)\) profinite for every oligomorphic group \(G\)?

\textbf{Unresolved qualifications.} A fully specified proposition requires the intended class of oligomorphic groups, the meaning of automorphism in \(\operatorname{Out}(G)\), and the topology on that outer automorphism group. If a quotient topology is intended, the topology on the automorphism group must also be specified. If only an abstract outer automorphism group is intended, the question must instead clarify whether it asks for the existence of a profinite group topology. None of these choices is settled by the supplied record.
\subsection{Short English statement}
The question asks whether every oligomorphic group has a profinite outer automorphism group. Oligomorphic means that its permutation action has finitely many orbits on tuples of every finite positive length. The intended automorphisms, topology, and possible closedness assumption require clarification.
\subsection{Sources}
\begin{itemize}
\item Local UnsolvedMath registry, supplied record with permanent ID \(801027\), exact source code \texttt{catalogue entry}; title: \emph{Open Problems in Mathematical Logic --- Problem 39}. The registry attributes the problem to Barmpalias, Gao et al., with proposed year \(2026\), and labels its status as open. This is a supplied status label, not an independently verified mathematical or literature assessment.
\item Supplied source attribution: Barmpalias, Gao et al., \emph{Open Problems in Mathematical Logic} (2026), Problem 39, \url{https://arxiv.org/abs/2608.26628}. The background records access on \texttt{2026-10-16} and extraction by \texttt{arxiv\_latex}; these provenance details are supplied and have not been independently verified here.
\end{itemize}
\end{document}
